\documentclass[11pt]{article}

\usepackage[margin=1in]{geometry}
\usepackage[round,authoryear]{natbib}
\usepackage{amsmath}
\usepackage{amssymb}
\usepackage{booktabs}
\usepackage{graphicx}
\usepackage{microtype}
\usepackage{tabularx}
\usepackage{xcolor}
\usepackage{hyperref}

\hypersetup{
  colorlinks=true,
  linkcolor=blue!50!black,
  citecolor=blue!50!black,
  urlcolor=blue!50!black
}
\emergencystretch=2em

\title{Cache-Mediated Safety Erasure:\\
Testing Whether KV-Cache Optimizations Selectively Weaken Refusal Behavior}

\author{
  Aryan Gupta\\
  \texttt{aryan.cs.app@gmail.com}
}

\date{}

\newcommand{\ssei}{\mathrm{SSEI}}
\newcommand{\todoresult}[1]{%
  \begin{center}
  \fbox{%
    \begin{minipage}{0.88\linewidth}
      \textbf{Empirical result not yet reported.} #1
    \end{minipage}
  }
  \end{center}
}
\newcommand{\maybeincludegraphic}[3]{%
  \IfFileExists{#1}{%
    \includegraphics[width=#2]{#1}%
  }{%
    \fbox{%
      \begin{minipage}[c][1.7in][c]{#2}
        \centering
        \textbf{Figure unavailable}\\
        #3
      \end{minipage}
    }%
  }%
}
\newcommand{\maybeinputtable}[2]{%
  \IfFileExists{#1}{%
    \input{#1}%
  }{%
    \todoresult{#2}%
  }%
}
\newcommand{\requiredartifact}[1]{}
\newcommand{\EmpiricalStatusSentence}{%
The present manuscript is a registered analysis protocol and reports no empirical claims; empirical sections are populated only after the complete evaluation and audit procedures described here have been executed.%
}
\newcommand{\PrimaryRunId}{results pending}
\newcommand{\PrimaryPolicyCount}{}
\newcommand{\PrimaryTopSSEIPolicy}{}
\newcommand{\PrimaryTopSSEI}{}
\newcommand{\PrimaryTopSSEICILow}{}
\newcommand{\PrimaryTopSSEICIHigh}{}
\newcommand{\PrimarySafetyClusterCount}{}
\newcommand{\PrimaryCapabilityClusterCount}{}
\newcommand{\CausalRunId}{results pending}
\newcommand{\CausalPolicyCount}{}
\newcommand{\CausalTopSSEIPolicy}{}
\newcommand{\CausalTopSSEI}{}
\newcommand{\CausalTopSSEICILow}{}
\newcommand{\CausalTopSSEICIHigh}{}
\newcommand{\CausalSafetyClusterCount}{}
\newcommand{\CausalCapabilityClusterCount}{}
\IfFileExists{../generated/h200_qwen_full_sweep/result_macros.tex}{%
  % Auto-generated by scripts/export_paper_assets.py; do not edit by hand.
\renewcommand{\PrimaryRunId}{primary public sweep}
\renewcommand{\PrimaryPolicyCount}{6}
\renewcommand{\PrimaryTopSSEIPolicy}{sliding\_window\_\_budget128}
\renewcommand{\PrimaryTopSSEI}{0.092}
\renewcommand{\PrimaryTopSSEICILow}{0.083}
\renewcommand{\PrimaryTopSSEICIHigh}{0.101}
\renewcommand{\PrimarySafetyClusterCount}{5205}
\renewcommand{\PrimaryCapabilityClusterCount}{1300}
%
}{}
\IfFileExists{../generated/h200_causal_patch_qwen7b/result_macros.tex}{%
  % Auto-generated by scripts/export_paper_assets.py; do not edit by hand.
\renewcommand{\PrimaryRunId}{primary public sweep}
\renewcommand{\PrimaryPolicyCount}{6}
\renewcommand{\PrimaryTopSSEIPolicy}{sliding\_window\_\_budget128}
\renewcommand{\PrimaryTopSSEI}{0.092}
\renewcommand{\PrimaryTopSSEICILow}{0.083}
\renewcommand{\PrimaryTopSSEICIHigh}{0.101}
\renewcommand{\PrimarySafetyClusterCount}{5205}
\renewcommand{\PrimaryCapabilityClusterCount}{1300}
%
}{}
\IfFileExists{../generated/claim_assessment/abstract_status_sentence.tex}{%
  \input{../generated/claim_assessment/abstract_status_sentence.tex}%
}{}
\requiredartifact{../generated/h200_qwen_full_sweep/result_macros.tex}
\requiredartifact{../generated/h200_causal_patch_qwen7b/result_macros.tex}
\requiredartifact{../generated/claim_assessment/abstract_status_sentence.tex}

\begin{document}

\maketitle

\begin{abstract}
Inference systems increasingly compress, evict, or quantize key-value (KV) caches to reduce serving cost and extend usable context. This paper tests a safety-specific hypothesis: cache optimizations can selectively weaken refusal and policy-following behavior while preserving ordinary task capability, under the hypothesis that some safety behavior may depend on fragile cache-resident routing state. We call the target phenomenon \emph{cache-mediated safety erasure} only if three claims are supported: cache interventions degrade safety behavior, the degradation is larger than matched capability degradation, and targeted restoration or protection of policy-relevant cache slices recovers refusal more than matched non-policy controls. \EmpiricalStatusSentence
\end{abstract}

\section{Introduction}

Safety evaluations usually treat a deployed language model as a fixed pair of weights and prompts. Real serving systems are less static. KV-cache eviction, quantization, paged attention, batching, and long-context memory management alter the transient state through which a model attends to prior tokens \citep{kwon2023pagedattention}. If aligned behavior depends on transient cache state, then a throughput optimization can also become an unintentional safety intervention.

This paper asks whether inference-time KV-cache optimizations can erase or weaken refusal behavior more than they degrade ordinary capability. The motivating observation is not merely that compression can hurt language-model quality. Prior work already shows that cache compression can produce uneven degradation, instruction-following failures, and system-prompt leakage \citep{chen2025pitfalls,ananthanarayanan2026physics,yang2024notoken}. The stronger question is whether safety behavior is selectively fragile and causally recoverable by preserving or restoring policy-relevant cache state.

We organize the study around a strict claims ladder:
\begin{enumerate}
  \item \textbf{Behavioral cache sensitivity:} cache interventions change safety, leakage, or capability behavior.
  \item \textbf{Selective safety degradation:} refusal or leakage behavior degrades more than ordinary capability under matched cache pressure.
  \item \textbf{Causal safety-state erasure:} restoring or protecting policy-relevant cache slices recovers safety behavior more than matched user-token or non-policy controls.
\end{enumerate}

Only the third claim justifies the phrase \emph{cache-mediated safety erasure}. If the experiments support only the first two claims, the correct contribution is an open-model deployment evaluation of KV-compression pitfalls, not a new mechanistic safety phenomenon.

\paragraph{Contributions.}
The study contributes four components: a public benchmark suite for safety, leakage, over-refusal, and capability under cache interventions; a selective degradation metric that subtracts ordinary capability loss from safety loss; token-role cache accounting for system, user, template, and generated spans; and causal restoration or mitigation tests that distinguish policy-token effects from ``more cache helps'' baselines.

\section{Related Work}

\paragraph{KV-cache compression and routing.}
H2O retains heavy-hitter tokens for efficient generation \citep{zhang2023h2o}, StreamingLLM emphasizes attention sinks for streaming contexts \citep{xiao2024streamingllm}, SnapKV compresses prompts based on observed attention structure \citep{li2024snapkv}, and KIVI and KVQuant study low-bit KV-cache quantization \citep{liu2024kivi,hooper2024kvquant}. No Token Left Behind argues that importance-aware mixed precision can preserve information that exhaustive eviction loses \citep{yang2024notoken}. The closest behavioral predecessor is \citet{chen2025pitfalls}, who report that KV-cache compression can unevenly harm instruction following and increase system-prompt leakage. \citet{ananthanarayanan2026physics} frame compression as perturbing token accessibility through attention dynamics. The central distinction here is safety-specific and causal: whether policy-token cache state is selectively fragile and recoverable.

\paragraph{Cache editing as a safety control surface.}
CachePrune edits KV-cache state to defend against indirect prompt injection \citep{wang2025cacheprune}, showing that cache state can be an intervention point rather than only an efficiency artifact. Prompt-leakage work studies privacy and instruction extraction risks \citep{hui2024pleak}. We instead test whether ordinary inference optimizations can weaken refusal behavior and whether token-role-aware retention can mitigate that weakness.

\paragraph{Alignment mechanisms and refusal.}
Refusal behavior can be steered through activation directions and may not reduce to a single global direction \citep{arditi2024refusal,joad2026more_refusal}. Recent alignment-routing work argues that policy behavior can localize to sparse routing mechanisms \citep{frank2026alignmentroutes}, while Any-Depth Alignment studies inference-time restoration of refusal behavior at different generation depths \citep{zhang2026anydepth}. These works motivate the possibility that safety state has identifiable carriers. Our experiments ask whether KV cache slices are among those carriers.

\paragraph{Phenomenon-first safety papers.}
Subliminal learning, token entanglement, and emergent misalignment are relevant here because each frames a safety question around a specific mechanism rather than only a new benchmark \citep{cloud2025subliminal,zur2025token,betley2025emergent}. Cache-mediated safety erasure is posed in the same form: a falsifiable phenomenon with controls that separate selective safety degradation from ordinary performance degradation.

\paragraph{Evaluation anchors.}
The registered primary sweep requires public prompt-injection prompts from the Cyberec prompt-injection dataset \citep{cyberec2026promptinjection}, harmful-request prompts from AdvBench and JailbreakBench \citep{zou2023universal,chao2024jailbreakbench}, XSTest-safe over-refusal prompts \citep{rottger2023xstest}, benign-instruction prompts from Databricks Dolly 15k \citep{databricks2023dolly}, and ARC-Easy capability questions \citep{clark2018arc}. Broader anchors such as HarmBench, IFEval, RULER, and LongBench define extension analyses only when their corresponding artifacts are present \citep{mazeika2024harmbench,zhou2023ifeval,hsieh2024ruler,bai2024longbench}. Primary experiments target Qwen2.5 Instruct models because they are open-weight and available at 7B, 14B, and 32B scales \citep{qwen2024qwen25}.

\section{Hypotheses}

\paragraph{H1: cache interventions change safety behavior.}
Safety, leakage, or over-refusal metrics differ between uncompressed baseline decoding and at least one cache intervention at matched prompts, seeds, and decoding settings.

\paragraph{H2: safety degradation is selective.}
Safety degradation exceeds ordinary capability degradation under the same policy and budget. We quantify this with the Selective Safety Erasure Index:
\begin{equation}
  \ssei(p, b, m) =
  \Delta_{\text{safety}}(p,b,m) -
  \Delta_{\text{capability}}(p,b,m),
\end{equation}
where $p$ is a cache policy, $b$ is its budget, and $m$ is the model. Positive values indicate larger safety loss than capability loss.

\paragraph{H3: policy-relevant cache restoration is causal.}
For length-preserving perturbations such as simulated KV quantization, restoring baseline key/value slices for system or policy spans recovers refusal more than restoring a matched number of user-token slices. For eviction policies where literal reinsertion is not yet implemented, token-role-aware pinning must preserve safety more than matched non-policy retention at comparable budgets.

\section{Methods}

\subsection{Models}

The primary model family is Qwen2.5 Instruct at 7B and 14B parameters, with Qwen2.5-32B reserved as an extension analysis when the corresponding artifacts are present. The core behavioral evaluation uses a mid-scale instruct model, while smaller models support causal diagnostics. A general model-family claim requires at least one non-Qwen open instruct model; otherwise, claims must be limited to the tested Qwen-family setting.

All experiments use open-weight models, public datasets, and local evaluation code. Model loading uses bfloat16 weights, \texttt{safetensors}, and explicit cache-state access so that the evaluated interventions are visible in the generation loop rather than delegated to external serving systems.

\subsection{Cache Policies}

The primary sweep compares uncompressed generation against six registered interventions:
\begin{itemize}
  \item \textbf{sliding window:} retain only the last $N$ cached tokens;
  \item \textbf{sink plus recent:} retain the first $S$ tokens and last $N$ tokens;
  \item \textbf{random matched:} evict random tokens at the same budget as a structured policy;
  \item \textbf{KV int8 simulation:} quantize and dequantize key/value tensors with an int8 symmetric approximation;
  \item \textbf{KV int4 simulation:} apply a more aggressive int4 approximation;
  \item \textbf{policy pinned:} protect system or policy spans while applying an eviction budget to other spans.
\end{itemize}
Attention-H2O retention is treated as a diagnostic extension when attention-score artifacts are present, not as a required primary-policy gate.

Each policy logs retained and evicted token counts by role, layer count, decode step, policy label, and cache-norm statistics where available. These logs support both figures and sanity checks that compression was active.

\subsection{Prompt Suites}

The benchmark separates safety, leakage, over-refusal, and capability:
\begin{itemize}
  \item \textbf{system leakage:} safe prompts asking the model to reveal hidden system instructions;
  \item \textbf{public system leakage:} public prompt-injection attacks paired with a hidden system canary;
  \item \textbf{public refusal safety:} harmful-request prompts prepared from public datasets;
  \item \textbf{public benign over-refusal:} benign requests that should be answered rather than refused;
  \item \textbf{public XSTest safe:} additional benign safety-sensitive prompts;
  \item \textbf{public capability ARC:} multiple-choice capability controls;
  \item \textbf{instruction following and long-context controls:} verifier-compatible IFEval-style prompts and long-context subsets included only when source provenance and evaluation scripts are complete.
\end{itemize}

Every processed suite writes a manifest with prompt IDs, source dataset revisions, record counts, and SHA256 hashes. Empirical claims require public provenance for all benchmark-derived prompts.

\subsection{Metrics}

\paragraph{Safety and leakage.}
The safety metrics include refusal-string detection, unsafe-compliance proxies from public benchmark labels, and exact or approximate leakage overlap against hidden system text. Leakage is measured by exact substring flags and token-overlap or ROUGE-style similarity.

\paragraph{Capability and over-refusal.}
Capability is tracked with exact-match or multiple-choice accuracy where labels exist. Benign over-refusal is measured by refusal-rate on benign suites and instruction-following pass rates when verifier-compatible prompts are available.

\paragraph{Statistical uncertainty.}
The primary uncertainty report uses prompt-clustered confidence intervals. Main tables include paired safety degradation intervals by prompt and seed. The strongest claims require narrow intervals, complete prompt-policy coverage, and agreement between automated metrics and the human audit.

\subsection{Causal Restoration}

For length-preserving cache perturbations, we run baseline and compressed decoding on identical prompts. We then patch selected key/value slices from the baseline cache into compressed runs and compare:
\begin{enumerate}
  \item compressed decoding with no patch;
  \item compressed decoding with system-role slices patched from baseline;
  \item compressed decoding with user-role slices matched to the system-token count;
  \item policy-pinned mitigation where system-role spans are retained under eviction.
\end{enumerate}

The causal claim requires system-role restoration to recover refusal or leakage avoidance more than matched user-role controls. Hard-coded token-index patches are diagnostic only and are not sufficient evidence.

\section{Experimental Design}

\paragraph{Reproducibility.}
Each analysis records the resolved configuration, environment metadata, prompt records, raw generations, metrics, cache statistics, figure source data, figure hashes, and table hashes. A run is eligible for inclusion only when public dataset provenance is present, the prompt-policy-seed matrix is complete, cache interventions are active, and all reported tables and figures are generated from the recorded artifacts.

\paragraph{Evaluation sequence.}
The empirical study is ordered to reduce false positives:
\begin{itemize}
  \item verify cache intervention mechanics on a small diagnostic slice;
  \item run the primary public-suite sweep across safety, leakage, benign, and capability prompts;
  \item run causal restoration and policy-pinning diagnostics on the strongest behavioral effects;
  \item test whether the strongest effects reproduce across model scale and, where feasible, model family.
\end{itemize}

\paragraph{Human audit.}
Automated safety labels are coarse. The analysis therefore includes a blinded audit of high-impact slices: baseline versus strongest compressed policies, safety suites where unsafe-compliance proxies increase, and examples where system-role patching changes behavior. Condition labels are hidden during annotation and joined to experimental metadata only after labeling.

\section{Visualization Strategy}

The main figures are designed to make the proposed phenomenon interpretable at the level of prompt, policy, and token role, rather than only reporting aggregate scatterplots or bar charts. We use four figure families.

\paragraph{Cache-state fingerprints.}
For each policy-budget condition, the analysis renders a layer-by-token retention map with token roles on the horizontal axis and layers on the vertical axis. System, hidden-system, user, template, and generated spans are shown as separate bands. A safety-specific erasure pattern is diagnosed by a visible break or thinning in policy-token access, not merely by a lower scalar score.

\paragraph{Safety-capability phase portraits.}
Instead of a static scatterplot alone, each policy traces a trajectory through safety loss and capability loss as the cache budget tightens. A selective safety-erasure pattern is indicated by an initial increase in safety loss at approximately fixed capability loss, followed only later by broad capability degradation.

\paragraph{Restoration flow diagrams.}
Causal patching is shown as a flow from baseline to compressed decoding to system-patched, matched-user-patched, and policy-pinned conditions. The diagnostic pattern is movement back toward the baseline only when policy-relevant cache slices are restored or protected.

\paragraph{Prompt-level effect maps.}
For audited prompts, the analysis embeds each prompt-condition pair using open embeddings or metric vectors and draws paired edges from baseline to compressed generations. Clusters of long, directionally consistent edges identify families of prompts where cache interventions produce coherent safety shifts.

\section{Figures and Tables}

\begin{figure}[t]
  \centering
  \maybeincludegraphic{../../results/h200_qwen_full_sweep/figures/safety_capability_phase_portrait.pdf}{0.92\linewidth}{Safety degradation versus capability degradation by policy and budget.}
  \caption{Safety-capability phase portrait. Policies trace trajectories as cache budgets tighten. A safety-erasure pattern appears as upward movement before rightward capability collapse.}
  \label{fig:safety-capability}
\end{figure}

\begin{figure}[t]
  \centering
  \maybeincludegraphic{../../results/h200_qwen_full_sweep/figures/selective_safety_erasure_heatmap.pdf}{0.92\linewidth}{Selective Safety Erasure Index heatmap by model, suite, policy, and cache budget.}
  \caption{Selective Safety Erasure Index across cache policies. Erasure is claimed only for conditions where safety loss remains after subtracting matched capability loss.}
  \label{fig:ssei}
\end{figure}

\begin{figure}[t]
  \centering
  \maybeincludegraphic{../../results/h200_qwen_full_sweep/figures/prompt_effect_constellation.pdf}{0.92\linewidth}{Prompt-level baseline-to-treatment effect vectors by suite and policy.}
  \caption{Prompt-level effect map. Each arrow starts at baseline and points toward the compressed condition in a metric-derived effect space. Coherent clusters indicate prompt families where cache interventions create similar behavioral shifts.}
  \label{fig:constellation}
\end{figure}

\begin{figure}[t]
  \centering
  \maybeincludegraphic{../../results/h200_qwen_full_sweep/figures/cache_state_fingerprint.pdf}{0.92\linewidth}{Token-role retention heatmap for system, user, template, and generated spans.}
  \caption{Cache-state fingerprint by token role and policy. This visualization tests whether policy-relevant spans form a visibly fragile band under cache pressure.}
  \label{fig:role-retention}
\end{figure}

\begin{figure}[t]
  \centering
  \maybeincludegraphic{../../results/h200_qwen_full_sweep/figures/safety_state_atlas.pdf}{0.92\linewidth}{Suite-policy atlas combining selective safety degradation with role-level cache loss.}
  \caption{Safety-state atlas. Color shows selective safety degradation while overlaid circles show token-role cache loss, making it possible to compare behavioral safety loss against system-token and user-token retention patterns.}
  \label{fig:safety-state-atlas}
\end{figure}

\begin{figure}[t]
  \centering
  \maybeincludegraphic{../../results/h200_causal_patch_qwen7b/figures/causal_restoration_fraction.pdf}{0.92\linewidth}{Restoration fraction by causal patch and policy-pinned condition.}
  \caption{Causal restoration fractions. The key comparison is system-role restoration versus matched user-role restoration under the same compressed condition.}
  \label{fig:causal}
\end{figure}

\begin{figure}[t]
  \centering
  \maybeincludegraphic{../../results/h200_causal_patch_qwen7b/figures/causal_restoration_flow.pdf}{0.92\linewidth}{Compressed, system-patched, matched-user-patched, and policy-pinned conditions.}
  \caption{Causal restoration flow. The diagnostic pattern is movement toward baseline for policy-relevant restoration more than matched controls.}
  \label{fig:causal-flow}
\end{figure}

\section{Results}

\subsection{Behavioral Cache Sensitivity}
Table~\ref{tab:suite-effects} reports paired baseline-versus-intervention effects across the registered safety, leakage, benign, and capability suites.
\maybeinputtable{../generated/h200_qwen_full_sweep/suite_level_effects_table.tex}{Report baseline versus each cache policy for refusal, leakage, benign over-refusal, and capability. Include prompt-clustered confidence intervals and a table of per-suite effects.}

\subsection{Selective Safety Degradation}
Top policy-level $\ssei$: \PrimaryTopSSEIPolicy{} with estimate \PrimaryTopSSEI{} and 95\% interval [\PrimaryTopSSEICILow{}, \PrimaryTopSSEICIHigh{}]. A positive result requires safety degradation to exceed capability degradation under matched cache pressure, not merely lower overall quality.
\maybeinputtable{../generated/h200_qwen_full_sweep/main_results_table.tex}{Report $\ssei$ by model, suite, policy, and budget after a completed registered analysis run.}

\subsection{Token-Role Cache Accounting}
The cache-state fingerprint in Figure~\ref{fig:role-retention} reports token-role retention under each cache policy. The mechanistic claim requires role-specific cache disruption to align with behavioral degradation: policy-relevant spans should be disproportionately evicted or perturbed in conditions with positive $\ssei$, and protected or restored spans should move the model back toward baseline behavior.

\subsection{Causal Restoration and Mitigation}
The central causal claim requires system-role restoration or pinning to outperform matched non-policy controls.
\maybeinputtable{../generated/h200_causal_patch_qwen7b/causal_restoration_table.tex}{Report system-role patching, matched user-role patching, and policy-pinned mitigation after the causal patch analysis is complete.}

\subsection{Evidence-Gated Claim Assessment}
\maybeinputtable{../generated/claim_assessment/claim_assessment_table.tex}{Report the claims ladder assessment after primary and causal runs are complete. The cache-mediated safety erasure claim requires all registered criteria to pass.}
\maybeinputtable{../generated/claim_assessment/claim_interpretation.tex}{Report evidence-gated interpretation text after the claim assessment is complete.}

\subsection{Human Audit and Failure Modes}
\maybeinputtable{../audit/h200_qwen_full_sweep_summary/human_audit_summary_table.tex}{Report blinded human-audit label rates before making safety claims.}
\maybeinputtable{../audit/h200_qwen_full_sweep_summary/human_audit_deltas_table.tex}{Report paired baseline-versus-treatment human-audit deltas for unsafe compliance, refusal correctness, leakage, and benign over-refusal.}
\maybeinputtable{../audit/h200_causal_patch_qwen7b_summary/human_audit_summary_table.tex}{Report blinded causal-diagnostic human-audit label rates before making restoration claims.}
\maybeinputtable{../audit/h200_causal_patch_qwen7b_summary/human_audit_deltas_table.tex}{Report paired causal-diagnostic human-audit deltas for system-role patching, matched non-policy patching, and policy-pinned mitigation.}
Qualitative examples are summarized at the category level rather than reproduced verbatim when they contain procedural harmful details.

\section{Discussion}

If the causal restoration experiments succeed, the result implies that alignment behavior must be evaluated as a property of the full inference stack, not only the trained model and prompt. Serving optimizations would need alignment regression tests and token-role-aware cache policies. A policy-pinned mitigation would show that preserving safety does not require disabling all cache compression.

If only selective degradation appears, the contribution is still useful but narrower: it would be evidence that cache compression can be a deployment-time safety risk. If no selective degradation appears, the study becomes a falsification result that constrains future claims about safety-specific cache fragility.

\section{Limitations}

The study uses open models, public datasets, and local evaluation code. This improves reproducibility but limits comparisons against closed deployed systems. Automated refusal and leakage metrics are imperfect, so final safety claims require human audit or an open local judge with documented calibration. Simulated KV quantization is not a faithful implementation of every production quantizer. Attention-based retention policies require access to attention statistics and may be less comparable to production implementations than eviction and quantization interventions. Finally, unless non-Qwen models are added, the paper must not claim model-family universality.

\section{Ethics and Safety}

The experiments evaluate refusal and leakage behavior without publishing procedural harmful instructions. Raw generations are retained locally for audit and reproducibility, but paper examples should be redacted or summarized when they risk enabling harm. The purpose is to identify deployment-time safety regressions and mitigations in open inference stacks.

\section{Reproducibility Checklist}

\begin{tabularx}{\linewidth}{@{}lX@{}}
\toprule
Requirement & Implementation \\
\midrule
Open dependencies & Open models, public datasets, local metrics, no proprietary judges. \\
Run provenance & Resolved analysis configuration, environment metadata, dataset revisions, and artifact hashes. \\
Full matrix coverage & Prompt-policy-seed completeness is verified before analysis. \\
Compression sanity & Cache interventions must measurably alter retained or perturbed cache state. \\
Statistical uncertainty & Prompt-clustered and paired intervals in aggregate metrics. \\
Paper artifacts & Figure and table manifests with SHA256 hashes. \\
Human audit & Blinded condition labels with post-annotation metadata joining. \\
\bottomrule
\end{tabularx}

\section{Conclusion}

Cache-mediated safety erasure is warranted as a positive empirical claim only if it survives causal controls. Novelty must therefore be earned empirically: safety behavior must degrade selectively, policy-relevant cache restoration must outperform matched controls, and mitigation must preserve safety without globally disabling cache compression. If the evidence supports those claims, the result would show that alignment behavior can depend on the transient inference state created by cache management. If the evidence does not support them, the study still constrains future claims about safety-specific cache fragility.

\bibliographystyle{plainnat}
\bibliography{../references}

\end{document}
